\section{An exact relative Hecke algebra for two separate cubic-twisted packets}
\label{sec:relative-cubic-hecke}

AC-G37 excludes putting two coincident-center scalar ladders
in one regular source parameter.  The following result instead
compares two separate eigensystems.  Its hypotheses concern
their Galois representations and a dense set of Frobenius
classes; it does not posit a new \texttt{NData} instance.

\begin{theorem}[AC-G38: the exact two-eigensystem congruence ideal]
\label{AC-G38}
Let $\mathcal O$ be the ring of integers of a finite extension
of $\mathbb Q_\ell$ containing a primitive cube root $\omega$.
Let $\Gamma$ be a profinite group, let
$\chi:\Gamma\to\mu_3$ be surjective and continuous, and let
$V:\Gamma\to\operatorname{GL}_n(\mathcal O)$ be continuous.
Put $H=\ker\chi$.  Suppose that over an algebraic closure
$\overline k$ of the residue field the residual
semisimplification has the form
\[
        \overline V^{\rm ss}=S\oplus\bigoplus_j\eta_j,
 \qquad \dim S>1,\quad \dim\eta_j=1,                 \tag{G38.1}
\]
where $S|_H$ is irreducible.  The constituent $S$ occurs
once in (G38.1); the scalar characters may repeat.
Let $\mathcal F$ be a set of conjugacy classes whose images
meet every conjugacy class in every finite quotient of
$\Gamma$.  Define $T\subset\mathcal O\times\mathcal O$
to be the $\mathcal O$-algebra generated by the diagonal
constants and the paired coefficients of
\[
             \det(X-V(g)),\qquad\det(X-\chi(g)V(g)),
                         \qquad g\in\mathcal F.
\]
Then
\[
 T=
 \begin{cases}
   \mathcal O\times\mathcal O,&\ell\ne3,\\[2pt]
   \{(a,b):b-a\in(\omega-1)\mathcal O\},&\ell=3.
 \end{cases}                                                \tag{G38.3}
\]
In particular, at $\ell=3$ the quotient
$(\mathcal O\times\mathcal O)/T$ is
$\mathcal O/(\omega-1)$ and has source-normalized length
\[
 \frac{\operatorname{length}_{\mathcal O}
                  ((\mathcal O\times\mathcal O)/T)}
      {\operatorname{length}_{\mathcal O}(\mathcal O/3)}
                           =\frac12.                 \tag{G38.4}
\]
\end{theorem}
\begin{proof}
\proofstep{1. Every two-eigensystem algebra is determined by one ideal.}
For any $\mathcal O$-subalgebra $T$ containing diagonal
constants, put
$I=\{b-a:(a,b)\in T\}\subset\mathcal O$.
It is an ideal: differences are additive and stable under
diagonal scalar multiplication.  Subtracting $(a,a)$ from
$(a,b)$ shows that $(0,b-a)\in T$, and therefore
\[
                         T=\{(a,a+i):a\in\mathcal O,\ i\in I\}.
                                                               \tag{G38.5}
\]
For the generated algebra in the statement, $I$ is the ideal
generated by the differences of its specified generators.
Indeed, the difference of a product satisfies
$bb'-aa'=b(b'-a')+a'(b-a)$, so forming products does not
enlarge that ideal beyond the ideal generated by those differences.

Every ideal of $\mathcal O$ is closed in its valuation topology.
The density hypothesis and continuity of characteristic
coefficients show that the coefficient differences for every
$g\in\Gamma$ belong to $I$.  Explicitly, modulo any power
of the maximal ideal the two representations and $\chi$
factor through one finite quotient.  A class from
$\mathcal F$ has the same image as $g$ up to conjugacy,
and hence the same two characteristic polynomials modulo
that power.  Taking the limit proves the claim.

\proofstep{2. A unit trace on every coset.}
For every fixed $g\in\Gamma$, we claim that some $h\in H$ has
\[
                 \operatorname{tr}V(gh)\in\mathcal O^\times.
                                                               \tag{G38.2}
\]
Let $W$ be the global semisimplification on the right
side of (G38.1), and restrict its action to $H$ when
forming the algebra below.
It contains $S|_H$ exactly once, since that constituent
has dimension greater than one and the other summands
are scalar.  The finite-dimensional algebra
$A=\operatorname{span}_{\overline k}\{W(h):h\in H\}$
acts faithfully and semisimply on $W$.
Its Jacobson radical annihilates each simple constituent,
so faithfulness makes the radical zero.  By the structure
of a finite-dimensional semisimple algebra over an
algebraically closed field, $A$ is the product of the
full matrix algebras for its inequivalent simple
constituents.  An exact primary locator is Etingof et al.,
\emph{Introduction to representation theory},
Theorem~3.2.2(ii), p.~46, and Theorem~3.5.4, p.~50,
\url{https://math.mit.edu/~etingof/reprbook.pdf}.
The density theorem and its preceding Proposition~3.1.4
and Corollary~3.2.1 were read with their full proofs here.
For the fixed $g$, consider the linear functional
$a\mapsto\operatorname{tr}(W(g)a)$ on $A$.
The global constituent $S$ is stable under $g$, and its
restriction to $H$ occurs once.  On its matrix factor
this functional is $a\mapsto\operatorname{tr}(S(g)a)$,
with $S(g)$ invertible.  Choose $a=S(g)^{-1}E_{11}$
in that factor and zero in the other factors.  The
functional takes value one.  It is therefore nonzero
on some spanning element $W(h)$, $h\in H$.
Trace is unchanged by semisimplification, proving (G38.2).  The multiplicity-one hypothesis
is essential here; a direct sum of three copies at
characteristic three would not give this conclusion.

\proofstep{3. Complete separation away from three.}
Suppose $\ell\ne3$.  Choose $g$ with $\chi(g)=\omega$
and choose $h\in H$ by (G38.2).  The trace-coefficient
difference at $gh$ is
$-(\omega-1)\operatorname{tr}V(gh)$.
Both factors are units in $\mathcal O$: a primitive cubic
root does not reduce to one away from characteristic three.
Step 1 puts this unit in $I$, so $I=\mathcal O$.
This argument makes no unramifiedness assumption at $269$
and applies also when $\ell=269$.

\proofstep{4. Exact depth at three.}
Now let $\ell=3$ and set $\delta=\omega-1$.
If $c_j(g)$ is the coefficient of $X^{n-j}$ in
$\det(X-V(g))$, the corresponding twisted coefficient
is $\chi(g)^j c_j(g)$.  Each difference is divisible
by $\delta$, so $I\subset\delta\mathcal O$.
Choose $g$ with $\chi(g)=\omega$ and $h\in H$
from (G38.2).  The difference of the trace coefficients
at $gh$ is $-(\omega-1)\operatorname{tr}V(gh)$,
a unit times $\delta$.  By Step 1 it belongs to $I$.  Hence
$I=\delta\mathcal O$.
Finally $\delta^2=-3\omega$, so
$v_{\mathcal O}(\delta)=v_{\mathcal O}(3)/2$.
Equation (G38.5) gives the quotient and (G38.4).
\end{proof}

\begin{corollary}[AC-G38-lattice: a bounded relative lattice loss]
\label{AC-G38-lattice}
Let a full $\mathcal O$-lattice $M\subset W_0\oplus W_1$
be stable under the algebra $T$ of (G38.3), acting through
its two scalar coordinates.  Let $P_i$ be the rational
coordinate projection and put
$M_i^+=P_iM$, $M_i^-=M\cap W_i$.
Then $M_i^+=M_i^-$ for $\ell\ne3$.  At $\ell=3$,
\[
 (\omega-1)M_i^+\subset M_i^-\subset M_i^+,
 \qquad
 \operatorname{lenUnit}_3(M_i^+/M_i^-)
                        \le\frac{\dim W_i}{2}.       \tag{G38.6}
\]
Thus, if these comparisons are identified with actual
source \texttt{etaJ} indices and their selected ranks
sum to at most \texttt{dZ0}, the additional bracket
in the normalization of (G28.6) is at most
$\tfrac12\log3$.
\end{corollary}
\begin{proof}
Away from three, $T$ contains both coordinate idempotents,
which therefore preserve $M$.  At three, it contains
$\delta P_0$ and $\delta P_1$, giving the inclusions.
The quotient is a quotient of
$M_i^+/\delta M_i^+$ and has length at most
$(\dim W_i)v_{\mathcal O}(\delta)$.
Dividing by $v_{\mathcal O}(3)$ gives (G38.6).
The conditional global assertion follows by the same
sum of residue degrees and ramification indices as
AC-G28-normalization, with the extra factor $1/2$.
It retains the explicit row identification and rank
hypotheses of that comparison.
\end{proof}

\paragraph{Exact relation to the supplied source.}
The residual shape (G38.1) is the good rank-$r$ constituent
plus the scalar ladder, with $r>1$ in \texttt{Params.Ranks}.
Its required absolute irreducibility on $\ker\chi_{269}$
is precisely the solvable-restriction conclusion used in
AC-G37-self; that conclusion is an explicit hypothesis of
the present application, not inferred from a polynomial
factorization alone.
The coset-trace argument makes an additional unramifiedness
assumption at $269$ unnecessary, even at coefficient three.
The only ramification condition used when obtaining the
source restriction conclusion is that $\chi_{269}$ be
unramified at the test primes.

The source's full Hecke coefficient comparison is
\texttt{IsHeckeGalois}, and its generators are specified
by \texttt{NData.T\_gen}.  The closure step is the same
one isolated in
\path{SumRepresentations.charpoly_clause_of_chebotarev}
and \texttt{T\_closed}, in
\path{ExistsInertiaCharpolyFactorOfTameType.lean}.
For genuine Frobenius elements the required density follows
from Chebotarev.  Restricting to places split over the maximal
real subfield preserves this density: in a finite Galois
extension of that real subfield, any element fixing the CM
field has trivial restriction to its quadratic quotient;
the primes realizing its conjugacy class therefore split
in that quotient.  Excluding finitely many places has no effect.

The theorem proves the exact algebra for the two specified
Galois eigensystems.  It does not assert that both eigensystems
have already been placed in one source theta lattice with
all \texttt{NStage} ties.  Away from three they belong to
different residual Hecke components; at three their exact
relative coefficient ideal is $(\omega-1)$.
Identifying the resulting relative lattice with the desired
\texttt{LinkN} index remains a new comparison theorem.
In particular, no upper bound on a twisted Selmer group
is deduced by reversing a congruence-index upper bound.
